% ============================================================
% RELATÓRIO FINAL PAC EXTENSIONISTA
% Católica de Santa Catarina
% ============================================================
 
\documentclass[12pt, a4paper]{article}
 
\usepackage[T1]{fontenc}
\usepackage[utf8]{inputenc}
\usepackage[brazil]{babel}
\usepackage{times}
\usepackage[top=3cm, bottom=2cm, left=3cm, right=2cm]{geometry}
\usepackage{setspace}
\usepackage{indentfirst}
\usepackage{titlesec}
\usepackage{graphicx}
\usepackage{xcolor}
\usepackage{fancyhdr}
\usepackage[hidelinks]{hyperref}
\usepackage{enumitem}
\usepackage{tocloft}
\usepackage{array}
\usepackage{tabularx}
\usepackage{booktabs}
\usepackage[alf]{abntex2cite}
\usepackage{longtable}
 
\setstretch{1.5}
\setlength{\parindent}{1.25cm}
\setlength{\parskip}{0pt}
 
\pagestyle{fancy}
\fancyhf{}
\fancyhead[R]{\thepage}
\renewcommand{\headrulewidth}{0pt}
 
\titleformat{\section}{\normalfont\bfseries\normalsize}{}{0pt}{}
\titleformat{\subsection}{\normalfont\bfseries\normalsize}{}{0pt}{}
\titleformat{\subsubsection}{\normalfont\bfseries\normalsize}{}{0pt}{}
\titlespacing*{\subsubsection}{0pt}{8pt}{4pt}
\titlespacing*{\section}{0pt}{12pt}{6pt}
\titlespacing*{\subsection}{0pt}{10pt}{4pt}
 
\renewcommand{\cftsecfont}{\bfseries}
\renewcommand{\cftsecpagefont}{\bfseries}
\setlength{\cftbeforesecskip}{2pt}
 
\begin{document}
 
% CAPA
\begin{titlepage}
  \centering
  \thispagestyle{fancy}
 
  {\bfseries CENTRO UNIVERSIT\'ARIO -- CAT\'OLICA DE SANTA CATARINA}
 
  \vspace{6pt}
  {\bfseries ENGENHARIA DE SOFTWARE}
 
  \vspace{24pt}
 
  {\bfseries
    Miguel Angelo Dufloth Filho \\[4pt]
    Miguel Angel Balladares Huertas \\[4pt]
    Leonardo Lot\'erio de Lima \\[4pt]
    Lucas Honorato dos Santos
  }
 
  \vfill
 
  {\bfseries RELAT\'ORIO FINAL}
 
  \vspace{6pt}
  {\bfseries PAC EXTENSIONISTA}
 
  \vspace{24pt}
 
  \includegraphics[width=8cm]{logo.png}
 
  \vfill
 
  {\bfseries JOINVILLE, 20 DE FEVEREIRO DE 2026}
\end{titlepage}
 
% FICHA DE APROVAÇÃO
\thispagestyle{fancy}
 
\begin{center}
  {\bfseries FICHA DE APROVAÇÃO}
\end{center}
 
\vspace{12pt}
 
\noindent RELAT\'ORIO do PAC Extensionista apresentado como requisito parcial de avalia\c{c}\~ao
na disciplina PAC Extensionista no curso de gradua\c{c}\~ao de Engenharia de Software do Centro
Universit\'ario Cat\'olica de Santa Catarina em Joinville, sob supervis\~ao e orienta\c{c}\~ao
do professor Claudinei Dias.
 
\vspace{24pt}
 
\noindent Joinville, 20 de fevereiro de 2026.
 
\vspace{36pt}
 
\begin{center}
  {\bfseries
    Miguel Angelo Dufloth Filho \\[4pt]
    Miguel Angel Balladares Huertas \\[4pt]
    Leonardo Lot\'erio de Lima \\[4pt]
    Lucas Honorato dos Santos
  }
 
  \vspace{16pt}
  Claudinei Dias
 
  \vspace{4pt}
  Professor respons\'avel
\end{center}
 
% SUMÁRIO
\newpage
\thispagestyle{fancy}
 
\begin{center}
  {\bfseries SUM\'ARIO}
\end{center}
 
\vspace{12pt}
 
\noindent\textbf{1 INTRODU\c{C}\~AO} \dotfill \pageref{sec:intro}
 
\vspace{4pt}
\noindent\textbf{2 DESCRI\c{C}\~AO DO P\'UBLICO BENEFICIADO PELAS A\c{C}\~OES DE EXTENS\~AO} \dotfill \pageref{sec:publico}
 
\vspace{4pt}
\noindent\textbf{3 OBJETIVOS} \dotfill \pageref{sec:objetivos}
 
\vspace{2pt}
\noindent\hspace{1em}3.1 OBJETIVO GERAL \dotfill \pageref{sec:obj_geral}
 
\vspace{2pt}
\noindent\hspace{1em}3.2 OBJETIVOS ESPEC\'IFICOS \dotfill \pageref{sec:obj_especificos}
 
\vspace{4pt}
\noindent\textbf{4 DESCRI\c{C}\~AO DAS PRINCIPAIS ATIVIDADES REALIZADAS} \dotfill \pageref{sec:atividades}
 
\vspace{4pt}
\noindent\textbf{5 AVALIA\c{C}\~AO DO PROJETO PELO P\'UBLICO BENEFICIADO} \dotfill \pageref{sec:avaliacao}
 
\vspace{4pt}
\noindent\textbf{6 CONSIDERA\c{C}\~OES FINAIS -- AUTOAVALIA\c{C}\~AO DO PAC EXTENSIONISTA} \dotfill \pageref{sec:consideracoes}
 
\vspace{4pt}
\noindent\textbf{REFER\^ENCIAS} \dotfill \pageref{sec:referencias}
 
\vspace{4pt}
\noindent\textbf{AP\^ENDICES} \dotfill \pageref{sec:apendices}
 
\vspace{4pt}
\noindent\textbf{ANEXOS} \dotfill \pageref{sec:anexos}
 
% 1 INTRODUÇÃO
\newpage
\phantomsection
\label{sec:intro}
\section*{1 INTRODU\c{C}\~AO}
 
O presente relat\'orio descreve o desenvolvimento de uma aplica\c{c}\~ao web denominada
\textbf{Torque Gest\~ao}, cujo objetivo \'e gerenciar servi\c{c}os e acompanhar processos de
mec\^anicas automotivas. A iniciativa surgiu da identifica\c{c}\~ao de uma lacuna tecnol\'ogica
no setor, onde estabelecimentos de pequeno e m\'edio porte frequentemente operam sem
ferramentas digitais adequadas para o controle de suas opera\c{c}\~oes.

A gest\~ao eficiente de uma oficina mec\^anica envolve m\'ultiplos processos simult\^aneos,
desde o atendimento ao cliente at\'e o controle das ordens de servi\c{c}o e o acompanhamento
do hist\'orico de cada ve\'iculo. A aus\^encia de sistemas informatizados nesse contexto pode
resultar em falhas de comunica\c{c}\~ao, perda de informa\c{c}\~oes e queda na qualidade do
atendimento prestado.

Nesse sentido, o PAC Extensionista prop\^os o desenvolvimento de uma solu\c{c}\~ao tecnol\'ogica
acess\'ivel e funcional, capaz de atender \`as demandas reais dessas oficinas. O sistema abrange
o gerenciamento interno de ordens de servi\c{c}o, o controle de clientes e ve\'iculos, e um
portal de autoatendimento para que o pr\'oprio cliente acompanhe o status do servi\c{c}o e
consulte o hist\'orico do seu ve\'iculo de forma aut\^onoma, contribuindo tanto para a melhoria
dos processos internos quanto para a transpar\^encia no atendimento prestado.
 
% 2 PÚBLICO BENEFICIADO
\newpage
\phantomsection
\label{sec:publico}
\section*{2 DESCRI\c{C}\~AO DO P\'UBLICO BENEFICIADO PELAS A\c{C}\~OES DE EXTENS\~AO}
 
O p\'ublico beneficiado pelas a\c{c}\~oes deste projeto extensionista s\~ao mec\^anicas automotivas
de pequeno e m\'edio porte. Trata-se de estabelecimentos que, em geral, possuem entre um e vinte
funcion\'arios, com opera\c{c}\~ao local ou regional, e que realizam servi\c{c}os de manuten\c{c}\~ao
preventiva e corretiva de ve\'iculos automotores.

O perfil desse p\'ublico \'e composto majoritariamente por empreendedores e gestores que conduzem
suas opera\c{c}\~oes de forma manual ou com recursos tecnol\'ogicos limitados, como planilhas e
anota\c{c}\~oes f\'isicas. Essa realidade imp\~oe desafios \`a organiza\c{c}\~ao dos servi\c{c}os,
ao controle financeiro e ao relacionamento com os clientes.

Al\'em dos gestores e mec\^anicos, o projeto beneficia diretamente os clientes finais das
oficinas --- propriet\'arios de ve\'iculos que passam a ter acesso a um portal de
autoatendimento para acompanhar o status dos servi\c{c}os em andamento e consultar o
hist\'orico completo de atendimentos de seus ve\'iculos, sem depender de contato telef\^onico
ou presencial com a oficina.

As a\c{c}\~oes do projeto foram direcionadas a oficinas situadas na regi\~ao de Joinville, Santa
Catarina, onde o setor automotivo representa uma parcela significativa da economia local, tornando
a proposta especialmente relevante para o contexto da comunidade atendida.
 
% 3 OBJETIVOS
\newpage
\phantomsection
\label{sec:objetivos}
\section*{3 OBJETIVOS}
 
\phantomsection
\label{sec:obj_geral}
\subsection*{3.1 Objetivo Geral}
 
Desenvolver uma aplica\c{c}\~ao web denominada Torque Gest\~ao para gerenciar servi\c{c}os,
ordens de servi\c{c}o e o relacionamento com clientes de mec\^anicas automotivas de pequeno e
m\'edio porte, proporcionando maior efici\^encia operacional e melhoria no atendimento prestado.
 
\phantomsection
\label{sec:obj_especificos}
\subsection*{3.2 Objetivos Espec\'ificos}
 
\begin{itemize}[leftmargin=2cm, itemsep=4pt]
  \item Implementar um m\'odulo de gerenciamento de ordens de servi\c{c}o, permitindo o
        registro, acompanhamento e encerramento de cada atendimento realizado na oficina.
  \item Desenvolver funcionalidades de gest\~ao de servi\c{c}os, possibilitando o cadastro
        e controle dos tipos de servi\c{c}os oferecidos pelo estabelecimento.
  \item Criar um sistema de relacionamento com clientes, centralizando o hist\'orico de
        atendimentos, dados dos ve\'iculos e informa\c{c}\~oes de contato de cada cliente.
  \item Disponibilizar ao cliente final um portal de autoatendimento para acompanhamento
        em tempo real do status dos servi\c{c}os em andamento e consulta ao hist\'orico
        de atendimentos de seus ve\'iculos, promovendo transpar\^encia e reduzindo a
        depend\^encia de contato direto com a oficina.
\end{itemize}
 
% 4 ATIVIDADES REALIZADAS
\newpage
\phantomsection
\label{sec:atividades}
\section*{4 DESCRI\c{C}\~AO DAS PRINCIPAIS ATIVIDADES REALIZADAS}
 
Para o desenvolvimento da aplica\c{c}\~ao Torque Gest\~ao, ser\~ao aplicadas as seguintes escolhas t\'ecnicas e de seguran\c{c}a, fundamentais para a estrutura\c{c}\~ao do projeto:
 
\subsection*{4.1 Stack Tecnol\'ogica}
 
\begin{itemize}[leftmargin=2cm, itemsep=4pt]
  \item \textbf{Linguagens de Programa\c{c}\~ao:} Python ser\'a utilizada para o desenvolvimento do back-end devido \`a sua sintaxe clara, alta produtividade e excelente ecossistema para a constru\c{c}\~ao de l\'ogicas de neg\'ocios e APIs. O JavaScript/TypeScript ser\'a adotado para o front-end por ser o padr\~ao da web para aplica\c{c}\~oes din\^amicas, onde o uso de tipagem est\'atica reduz erros em tempo de execu\c{c}\~ao e facilita a manuten\c{c}\~ao.
  \item \textbf{Frameworks e Bibliotecas:} Ser\'a utilizado o FastAPI no back-end para a cria\c{c}\~ao \'agil e robusta de APIs RESTful, garantindo alta performance e seguran\c{c}a. O React ser\'a adotado no front-end para construir a interface de usu\'ario de forma componentizada, o que facilita a cria\c{c}\~ao de pain\'eis interativos (dashboards) para acompanhamento em tempo real das ordens de servi\c{c}o.
  \item \textbf{Banco de Dados:} O PostgreSQL ser\'a empregado como sistema gerenciador de banco de dados relacional (SGBDR), garantindo alta integridade, seguran\c{c}a e escalabilidade para o armazenamento das informa\c{c}\~oes sens\'iveis de clientes, ve\'iculos e hist\'oricos de ordens de servi\c{c}o.
  \item \textbf{Ferramentas de Desenvolvimento e Gest\~ao:} O desenvolvimento ocorrer\'a com o uso da IDE Visual Studio Code. O versionamento de c\'odigo ser\'a realizado utilizando Git, com reposit\'orios no GitHub/GitLab. Para conteineriza\c{c}\~ao, ser\'a adotado o Docker, assegurando a paridade entre os ambientes de desenvolvimento e produ\c{c}\~ao. A automa\c{c}\~ao de testes e deploys (CI/CD) ser\'a estruturada via GitHub Actions.
  \item \textbf{Hospedagem:} A infraestrutura da aplica\c{c}\~ao ser\'a baseada em nuvem, utilizando a plataforma Render para a hospedagem do back-end (FastAPI) e do banco de dados (PostgreSQL), e o Vercel para a entrega otimizada do front-end (React), assegurando alta disponibilidade e facilidade de escala dentro do contexto acad\^emico do projeto.
  \item \textbf{Licenciamento:} O projeto priorizar\'a tecnologias de c\'odigo aberto (Open Source), incorporando ferramentas sob a MIT License (React e FastAPI) e Apache License 2.0 (Docker).
\end{itemize}
 
\subsection*{4.2 Considera\c{c}\~oes de Seguran\c{c}a}
 
\begin{itemize}[leftmargin=2cm, itemsep=4pt]
  \item \textbf{Riscos Identificados:} Est\~ao mapeadas amea\c{c}as como o vazamento de dados pessoais e veiculares de clientes, inje\c{c}\~ao de c\'odigo (SQL Injection) atrav\'es de formul\'arios de cadastro, e falhas de autentica\c{c}\~ao que poderiam permitir acessos indevidos a m\'odulos administrativos restritos.
  \item \textbf{Medidas de Mitiga\c{c}\~ao:} Para minimizar os riscos, ser\'a aplicada a criptografia de senhas utilizando Bcrypt e tr\'afego de dados sob o protocolo HTTPS (TLS/SSL). Ser\'a utilizado um ORM (Object-Relational Mapping) para sanitiza\c{c}\~ao autom\'atica das entradas no banco de dados. O controle de acesso ser\'a implementado via autentica\c{c}\~ao baseada em tokens (JWT) em conjunto com o modelo RBAC (\textit{Role-Based Access Control}).
  \item \textbf{Normas e Boas Pr\'aticas Seguidas:} O projeto basear-se-\'a nas diretrizes do OWASP Top 10 para a preven\c{c}\~ao das vulnerabilidades web mais cr\'iticas e observar\'a os preceitos da Lei Geral de Prote\c{c}\~ao de Dados (LGPD) no que tange \`a minimiza\c{c}\~ao da coleta e \`a necessidade de consentimento.
  \item \textbf{Responsabilidade \'Etica:} A privacidade ser\'a tratada como pilar no tratamento dos dados de clientes, garantindo total transpar\^encia, confidencialidade e a aus\^encia de compartilhamento de informa\c{c}\~oes com terceiros. A abordagem assegurar\'a a titularidade da informa\c{c}\~ao do consumidor e o uso exclusivo para a presta\c{c}\~ao do servi\c{c}o automotivo.
\end{itemize}
 
\subsection*{4.3 Plano de Gerenciamento dos Riscos}
 
O plano de gerenciamento dos riscos descreve como os riscos do projeto Torque Gest\~ao ser\~ao
identificados, avaliados, monitorados e controlados ao longo do seu ciclo de vida, servindo
como base estrat\'egica para mitigar incertezas e aumentar as chances de sucesso.
 
\subsubsection*{4.3.1 Metodologia}
 
A abordagem adotada para o gerenciamento dos riscos combina an\'alise qualitativa e quantitativa.
Os riscos ser\~ao identificados por meio de sess\~oes de \textit{brainstorming} entre os membros
da equipe, revis\~ao de li\c{c}\~oes aprendidas de projetos similares e an\'alise dos requisitos
t\'ecnicos do sistema. Cada risco identificado ser\'a registrado, classificado por categoria,
avaliado quanto \`a probabilidade e ao impacto, e ter\'a uma resposta planejada associada.
 
\subsubsection*{4.3.2 Fun\c{c}\~oes e Responsabilidades}
 
\vspace{6pt}
\noindent
\renewcommand{\arraystretch}{1.4}
\begin{tabular}{|p{5cm}|p{9cm}|}
\hline
\textbf{Respons\'avel} & \textbf{Atribui\c{c}\~oes} \\
\hline
Equipe de desenvolvimento & Identificar e reportar novos riscos t\'ecnicos durante o desenvolvimento \\
\hline
L\'ider do projeto & Avaliar, priorizar e monitorar os riscos registrados \\
\hline
Todos os integrantes & Participar das revis\~oes peri\'odicas e atualizar o status dos riscos \\
\hline
Professor orientador & Validar as estrat\'egias de resposta e acompanhar a evolu\c{c}\~ao do plano \\
\hline
\end{tabular}
 
\subsubsection*{4.3.3 Or\c{c}amenta\c{c}\~ao}
 
Por se tratar de um projeto acad\^emico com infraestrutura baseada em servi\c{c}os gratuitos ou
de baixo custo (planos \textit{free tier} de AWS, Render e Vercel), os custos diretos de
mitiga\c{c}\~ao de riscos s\~ao m\'inimos. Os recursos necess\'arios envolvem principalmente
tempo de desenvolvimento para implementa\c{c}\~ao das medidas de seguran\c{c}a, como
criptografia, testes de vulnerabilidade e configura\c{c}\~ao de autentica\c{c}\~ao.
 
\subsubsection*{4.3.4 Tempos}
 
As atividades de gerenciamento de riscos ser\~ao realizadas conforme o cronograma abaixo:
 
\vspace{6pt}
\noindent
\renewcommand{\arraystretch}{1.4}
\begin{tabular}{|p{6cm}|p{8cm}|}
\hline
\textbf{Atividade} & \textbf{Frequ\^encia / Momento} \\
\hline
Identifica\c{c}\~ao de novos riscos & A cada sprint ou entrega de m\'odulo \\
\hline
Revis\~ao da matriz de riscos & Quinzenalmente \\
\hline
Atualiza\c{c}\~ao do status dos riscos & Ap\'os cada reuni\~ao de equipe \\
\hline
Auditoria final de riscos & Antes da entrega do produto final \\
\hline
\end{tabular}
 
\subsubsection*{4.3.5 Categorias de Risco}
 
Os riscos do projeto foram agrupados nas seguintes categorias:
 
\begin{itemize}[leftmargin=2cm, itemsep=4pt]
  \item \textbf{T\'ecnicos:} falhas de autentica\c{c}\~ao, vulnerabilidades de c\'odigo, indisponibilidade de infraestrutura e incompatibilidade entre tecnologias.
  \item \textbf{Segurança e Privacidade:} vazamento de dados, SQL Injection e viola\c{c}\~ao da LGPD.
  \item \textbf{Organizacionais:} atrasos no cronograma, sobrecarga da equipe e mudan\c{c}as de escopo.
  \item \textbf{Externos:} descontinuidade de servi\c{c}os terceiros (APIs, plataformas de hospedagem) e altera\c{c}\~oes em requisitos legais.
\end{itemize}
 
\subsubsection*{4.3.6 Defini\c{c}\~oes de Probabilidade e Impacto}
 
\vspace{6pt}
\noindent
\renewcommand{\arraystretch}{1.4}
\begin{tabular}{|c|p{5.5cm}|p{5.5cm}|}
\hline
\textbf{N\'ivel} & \textbf{Probabilidade} & \textbf{Impacto} \\
\hline
Baixo & Pouco prov\'avel de ocorrer no ciclo do projeto & Efeito m\'inimo, sem comprometer entregas \\
\hline
M\'edio & Poss\'ivel ocorr\^encia durante o desenvolvimento & Atraso ou retrabalho pontual em m\'odulos \\
\hline
Alto & Prov\'avel ou j\'a ocorreu em projetos similares & Compromete funcionalidades cr\'iticas ou a entrega final \\
\hline
\end{tabular}
 
\subsubsection*{4.3.7 Matriz de Probabilidade e Impacto}
 
A matriz a seguir combina probabilidade e impacto para classificar o n\'ivel de cada risco,
orientando a prioriza\c{c}\~ao das a\c{c}\~oes de resposta. Riscos classificados como
\textbf{Alto} exigem a\c{c}\~ao imediata; \textbf{M\'edio} requerem monitoramento cont\'inuo;
\textbf{Baixo} s\~ao aceit\'aveis com controle peri\'odico.
 
\vspace{6pt}
\noindent
\renewcommand{\arraystretch}{1.4}
\begin{tabular}{|p{4.2cm}|c|c|c|p{3.8cm}|}
\hline
\textbf{Risco} & \textbf{Categoria} & \textbf{Prob.} & \textbf{Impacto} & \textbf{N\'ivel} \\
\hline
Vazamento de dados de clientes & Seguran\c{c}a & Baixa & Alto & M\'edio \\
\hline
SQL Injection em formul\'arios & Seguran\c{c}a & M\'edia & Alto & Alto \\
\hline
Falha de autentica\c{c}\~ao & T\'ecnico & Baixa & Alto & M\'edio \\
\hline
Indisponibilidade da infraestrutura & T\'ecnico & Baixa & M\'edio & Baixo \\
\hline
Viola\c{c}\~ao da LGPD & Privacidade & Baixa & Alto & M\'edio \\
\hline
Atraso no cronograma & Organizacional & M\'edia & M\'edio & M\'edio \\
\hline
Descontinuidade de servi\c{c}o externo & Externo & Baixa & M\'edio & Baixo \\
\hline
\end{tabular}
 
\subsubsection*{4.3.8 Toler\^ancias das Partes Interessadas}
 
As principais partes interessadas do projeto s\~ao a equipe de desenvolvimento, o professor
orientador e as oficinas mec\^anicas beneficiadas. A toler\^ancia ao risco \'e considerada
\textbf{baixa} para riscos de seguran\c{c}a e privacidade, dado o car\'ater sens\'ivel dos
dados de clientes e ve\'iculos. Para riscos organizacionais, como atrasos pontuais no
cronograma acad\^emico, a toler\^ancia \'e \textbf{moderada}, desde que as entregas principais
sejam mantidas.
 
\subsubsection*{4.3.9 Formatos de Relat\'orios}
 
Os riscos ser\~ao documentados e comunicados por meio de: atualiza\c{c}\~oes na matriz de
riscos ap\'os cada revis\~ao; relat\'orios de status compartilhados com o professor orientador;
e registro das ocorr\^encias e respostas aplicadas no reposit\'orio do projeto (GitHub/GitLab),
garantindo rastreabilidade e transpar\^encia ao longo do desenvolvimento.
 
\subsubsection*{4.3.10 Forma de Acompanhamento}
 
O monitoramento dos riscos ser\'a realizado por meio de reuni\~oes quinzenais da equipe,
inspe\c{c}\~oes de c\'odigo com foco em seguran\c{c}a (OWASP Top 10), uso de ferramentas
de an\'alise est\'atica e testes de integra\c{c}\~ao cont\'inua via GitHub Actions. Novos
riscos identificados durante o desenvolvimento ser\~ao incorporados imediatamente \`a matriz,
e riscos encerrados ser\~ao registrados como li\c{c}\~oes aprendidas para projetos futuros.

% ============================================================
%  SEÇÃO 4.4 — PROTOTIPAÇÃO E INTEGRAÇÕES
% ============================================================

\subsection*{4.4 Prototipa\c{c}\~ao de Alta Fidelidade e Defini\c{c}\~ao de Integrações}

O processo de design do Torque Gest\~ao foi conduzido com base em um prot\'otipo de alta
fidelidade, desenvolvido com o objetivo de validar os fluxos de navega\c{c}\~ao, a
consist\^encia visual e a usabilidade das telas antes do in\'icio da implementa\c{c}\~ao.
A prototipa\c{c}\~ao antecede o desenvolvimento como estrat\'egia para reduzir o risco de
retrabalho e garantir o alinhamento entre a equipe e os requisitos especificados.

\subsubsection*{4.4.1 Ferramenta e Design System}

O prot\'otipo foi constru\'ido no \textbf{Figma}, ferramenta de design colaborativo adotada
pelo projeto para a cria\c{c}\~ao de todos os \textit{wireframes} e telas de alta fidelidade.
O design system aplicado em todas as telas segue as seguintes defini\c{c}\~oes:

\begin{itemize}[leftmargin=2cm, itemsep=4pt]
  \item \textbf{Cor prim\'aria:} \texttt{\#1B2B4B} (azul escuro institucional)
  \item \textbf{Cor de destaque:} \texttt{\#F0A500} (âmbar)
  \item \textbf{Tipografia:} Inter
  \item \textbf{Princ\'ipios de interface:} elementos de grande dimens\~ao e alto contraste,
        derivados da observa\c{c}\~ao direta do ambiente operacional da oficina (ver
        Subse\c{c}\~ao~\ref{subsec:elicitacao})
\end{itemize}

O prot\'otipo cobre os seguintes fluxos principais, organizados em conformidade com os
requisitos funcionais especificados na Se\c{c}\~ao~4.5: autentica\c{c}\~ao e dashboards
por perfil (RF06); gerenciamento de clientes e ve\'iculos (RF01, RF02); ciclo completo
de ordens de servi\c{c}o (RF03, RF04); cat\'alogo de servi\c{c}os e pe\c{c}as (RF05);
gest\~ao de usu\'arios e configura\c{c}\~oes; e portal do cliente (RF07).

O acesso ao prot\'otipo interativo est\'a dispon\'ivel no seguinte endere\c{c}o:

\begin{center}
  \textbf{[INSERIR LINK DO FIGMA]}
\end{center}

\subsubsection*{4.4.2 Integrações Previstas}

O Torque Gest\~ao foi projetado com uma arquitetura aberta que permite integrações futuras
por meio da API RESTful exposta pelo back-end (RNF06). As integrações previstas para
vers\~oes subsequentes ao MVP s\~ao:

\begin{itemize}[leftmargin=2cm, itemsep=4pt]
  \item \textbf{Fornecedores de pe\c{c}as:} consulta de disponibilidade e pre\c{c}os via
        API, agilizando a composi\c{c}\~ao de itens na ordem de servi\c{c}o.
  \item \textbf{Notifica\c{c}\~oes externas:} envio de atualiza\c{c}\~oes de status ao
        cliente por e-mail ou SMS, complementando o portal de autoatendimento (RF07).
  \item \textbf{Sistemas de gest\~ao cont\'abil:} integrações com plataformas
        especializadas (e.g., Conta Azul) para repasse das informa\c{c}\~oes operacionais
        ao controle fiscal e tribut\'ario, mantendo o escopo do Torque Gest\~ao restrito
        \`a gest\~ao da oficina.
\end{itemize}

% ============================================================
%  SEÇÃO 4.5 — ENGENHARIA DE REQUISITOS
% ============================================================

\section{4.5 Engenharia de Requisitos}

A Engenharia de Requisitos (ER) constitui uma das fases mais críticas e
determinantes para o sucesso do \textit{software} Torque Gestão, atuando como
interface fundamental entre as necessidades de negócio — frequentemente
informais e subjetivas — dos gestores de oficinas mecânicas e a implementação
técnica rigorosa. Conforme preconizado por \citeonline{sommerville2011}, a ER
não é apenas uma etapa de escrita, mas um processo cíclico de descoberta,
análise, documentação e verificação dos serviços e restrições que o sistema
deve apresentar.

Neste projeto, a ER foi tratada como pilar estratégico para mitigar o risco de
desenvolvimento de funcionalidades inúteis ou de difícil operação. O objetivo
central foi transformar a ausência de ferramentas de gestão informatizada
identificada no setor automotivo em uma especificação técnica que garantisse
eficiência operacional, segurança jurídica — em especial a conformidade com a
Lei Geral de Proteção de Dados (LGPD) — e uma curva de aprendizado reduzida
para os profissionais do setor.

\citeonline{sommerville2011} distingue três categorias de requisitos que
orientaram toda a estruturação desta seção: \textbf{requisitos funcionais},
que descrevem os serviços e comportamentos esperados do sistema;
\textbf{requisitos não funcionais}, que especificam restrições de qualidade e
propriedades globais do sistema; e \textbf{requisitos de domínio}, oriundos do
contexto de aplicação, como regulamentações legais e normas setoriais. As
tabelas de requisitos apresentadas nas subseções seguintes compõem o núcleo do
Documento de Requisitos do sistema, elaborado conforme as diretrizes do padrão
IEEE~830.

% ------------------------------------------------------------
\subsection{4.5.1 Elicitação e Análise de Requisitos}
\label{subsec:elicitacao}
% ------------------------------------------------------------

A fase de elicitação buscou compreender profundamente o domínio do problema: o
déficit de maturidade digital e a gestão predominantemente analógica em
oficinas mecânicas de pequeno e médio porte na região de Joinville, Santa
Catarina. Dado que o setor automotivo representa uma parcela significativa da
economia local, a precisão nesta fase foi essencial para garantir que os
requisitos levantados refletissem necessidades reais e não suposições da equipe
de desenvolvimento.

Foram adotadas as seguintes técnicas complementares para assegurar uma visão
holística do problema:

\subsubsection{4.5.1.1 Entrevistas Estruturadas e Semiestruturadas}

Foram realizadas reuniões com proprietários e administradores de oficinas. O
foco não se limitou a identificar \textit{"o que o sistema deve fazer"}, mas a
compreender as principais dificuldades operacionais recorrentes no processo de
atendimento. Identificou-se que a perda de prazos de manutenção preventiva e a
dificuldade em localizar o histórico de veículos com atendimentos anteriores
eram os principais gargalos financeiros enfrentados pelos estabelecimentos.

\subsubsection{4.5.1.2 Observação Direta (\textit{Etnografia})}

A equipe analisou o fluxo de trabalho real no ambiente operacional da oficina.
Observou-se que os mecânicos frequentemente operam em condições que dificultam
a interação com dispositivos convencionais — mãos contaminadas por graxa e
posições ergonomicamente desfavoráveis ao uso de teclados. Essa observação
evidenciou a necessidade de interfaces com elementos visuais de grande dimensão,
alto contraste cromático e suporte a dispositivos móveis (\textit{tablets}),
reduzindo a dependência de registros físicos que podem se extraviar ou
tornar-se ilegíveis.

\subsubsection{4.5.1.3 Análise de Cenários e \textit{User Stories}}

Foram elaboradas narrativas de uso para validar a lógica dos fluxos do sistema
e identificar requisitos implícitos não mencionados nas entrevistas. Um exemplo
representativo é apresentado a seguir:

\begin{quote}
\textit{"Um cliente retorna após seis meses com um problema recorrente no
veículo. O mecânico responsável pelo atendimento precisa recuperar o diagnóstico
anterior registrado no sistema para comparar as falhas e evitar retrabalho."}
\end{quote}

\noindent\textbf{Critério de aceitação:} o sistema deve exibir o último
diagnóstico registrado para o veículo em até três interações a partir da tela
de consulta de Ordens de Serviço, sem necessidade de buscas manuais em
registros físicos.

Esse exercício permitiu identificar requisitos implícitos de rastreabilidade e
agilidade de consulta que não haviam emergido nas entrevistas estruturadas.
Ao final da fase de elicitação, compreendeu-se que a agilidade na visualização
de status e a integridade dos dados são mais valorizadas pelos usuários do que
um grande volume de funcionalidades complexas.

% ------------------------------------------------------------
\subsection{4.5.2 Especificação de Requisitos}
\label{subsec:especificacao}
% ------------------------------------------------------------

A especificação define detalhadamente as funções e as restrições operacionais
do sistema. Seguindo o modelo proposto por \citeonline{sommerville2011}, os
requisitos são organizados em dois níveis de abstração: \textbf{Requisitos de
Usuário}, expressos em linguagem natural e destinados aos
\textit{stakeholders} não técnicos, e \textbf{Requisitos de Sistema},
descrições técnicas detalhadas destinadas à equipe de desenvolvimento. As
tabelas a seguir representam os Requisitos de Sistema do Torque Gestão, já
refinados a partir das declarações de usuário coletadas na fase de elicitação.

% ------------------------------------------------------------
\subsubsection{4.5.2.1 Requisitos Funcionais}
% ------------------------------------------------------------

Os requisitos funcionais descrevem os serviços específicos que o
\textit{software} deve fornecer, o modo como o sistema deve reagir a entradas
específicas e como deve se comportar em determinadas situações
\cite{sommerville2011}. Os requisitos identificados para o Torque Gestão são
apresentados na Tabela~\ref{tab:rf}.

%Tabela 1 - Requisitos funcionais

\begin{longtable}{|p{1.2cm}|p{3.5cm}|p{9.5cm}|}
\caption{Requisitos Funcionais do Torque Gestão}
\label{tab:rf} \\

\hline
\textbf{ID} & \textbf{Título} & \textbf{Descrição e Implicações} \\
\hline
\endfirsthead

\multicolumn{3}{r}{\footnotesize\textit{(continuação)}} \\
\hline
\textbf{ID} & \textbf{Título} & \textbf{Descrição e Implicações} \\
\hline
\endhead

\hline
\multicolumn{3}{r}{\footnotesize\textit{(continua na próxima página)}} \\
\endfoot

\hline
\endlastfoot

RF01 & Gerenciamento de Clientes &
O sistema deve permitir o cadastro completo, a edição e a consulta de clientes,
com validação de CPF/CNPJ e mecanismo de busca por nome ou documento para
prevenir a duplicidade de registros. \\
\hline
RF02 & Cadastro de Veículos &
O sistema deve vincular múltiplos veículos a um único proprietário, registrando
Placa (formato Mercosul --- \textbf{requisito de domínio}, ver
Subseção~\ref{subsubsec:dominio}), Modelo, Marca, Ano e quilometragem atual
para subsidiar previsões de manutenção preventiva. \\
\hline
RF03 & Emissão de Ordens de Serviço (OS) &
O sistema deve permitir a abertura de OS detalhadas, separando obrigatoriamente
os itens de \textit{"Mão de Obra"} dos itens de \textit{"Peças Aplicadas"},
possibilitando a geração de orçamentos prévios antes da aprovação final pelo
cliente. Esta separação implica um modelo relacional com tabela de itens
discriminados por tipo, justificando o uso do PostgreSQL como SGBD. \\
\hline
RF04 & Acompanhamento de Status &
O sistema deve implementar um fluxo de estados para a OS: \textit{"Aguardando
Diagnóstico"}, \textit{"Em Execução"}, \textit{"Aguardando Peças"},
\textit{"Finalizada"} e \textit{"Entregue"}, possibilitando o controle de
produtividade da equipe e a comunicação proativa com o cliente. \\
\hline
RF05 & Catálogo de Serviços e Peças &
O sistema deve manter um cadastro de serviços (e.g., alinhamento, troca de
óleo) e peças, com descrição e referência de preço operacional para uso interno
da oficina, agilizando a composição de orçamentos durante o atendimento.
O controle fiscal e tributário é responsabilidade de sistemas especializados
externos (e.g., sistemas de gestão contábil), fora do escopo deste projeto. \\
\hline
RF06 & Autenticação e Controle de Acesso &
O sistema deve validar credenciais de acesso e aplicar o Controle de Acesso
Baseado em Funções (\textit{Role-Based Access Control} --- RBAC), por meio de
tokens JWT. São definidos três perfis: \textbf{Administrador}, com acesso
completo ao sistema, incluindo configurações e relatórios operacionais;
\textbf{Mecânico}, com acesso restrito às OS sob sua responsabilidade;
e \textbf{Cliente}, com acesso exclusivo ao portal de acompanhamento de seus
próprios veículos e serviços (ver RF07). \\
\hline
RF07 & Portal do Cliente &
O sistema deve disponibilizar um portal de acesso \textit{self-service} para o
cliente final, acessível via credenciais próprias. Por meio deste portal, o
cliente poderá: (i) consultar o histórico completo de OS de seus veículos
vinculados; (ii) acompanhar em tempo real o status da OS em andamento,
refletindo os estados definidos em RF04; e (iii) visualizar o orçamento aprovado
e os serviços executados. O portal não expõe dados de outros clientes, sendo o
isolamento de dados garantido pela camada de autorização RBAC. \\

\end{longtable}

% ------------------------------------------------------------
\subsubsection{4.5.2.2 Requisitos Não Funcionais}
% ------------------------------------------------------------

Os requisitos não funcionais especificam critérios de qualidade e restrições
técnicas que o sistema deve observar, não descrevendo o que o sistema faz, mas
como o faz \cite{sommerville2011}. Eles frequentemente se aplicam ao sistema
como um todo e seu descumprimento pode inviabilizar a adoção do
\textit{software} pelo público-alvo. Os requisitos não funcionais do Torque
Gestão são apresentados na Tabela~\ref{tab:rnf}.

%Tabela 2 - Requisitos não funcionais

\begin{longtable}{|p{1.2cm}|p{3.5cm}|p{9.5cm}|}
\caption{Requisitos Não Funcionais do Torque Gestão}
\label{tab:rnf} \\

\hline
\textbf{ID} & \textbf{Título} & \textbf{Descrição e Implicações} \\
\hline
\endfirsthead

\multicolumn{3}{r}{\footnotesize\textit{(continuação)}} \\
\hline
\textbf{ID} & \textbf{Título} & \textbf{Descrição e Implicações} \\
\hline
\endhead

\hline
\multicolumn{3}{r}{\footnotesize\textit{(continua na próxima página)}} \\
\endfoot

\hline
\endlastfoot

RNF01 & Segurança de Dados &
Todas as credenciais de acesso devem ser armazenadas com \textit{hash}
criptográfico Bcrypt no PostgreSQL, impedindo a leitura de senhas em texto
claro em caso de acesso não autorizado ao banco de dados. O tráfego entre
cliente e servidor deve ser protegido pelo protocolo HTTPS (TLS/SSL). \\
\hline
RNF02 & Desempenho &
O sistema deve ser otimizado de modo que consultas de histórico de OS — em
condições de uso típico com até 500 registros por veículo — não ultrapassem
2~segundos de tempo de resposta, mensuráveis por ferramentas de monitoramento
de latência de API. Tempos superiores comprometem o fluxo de atendimento no
ambiente operacional da oficina. \\
\hline
RNF03 & Disponibilidade &
A aplicação será hospedada em infraestrutura de nuvem, utilizando a plataforma
Render para o \textit{back-end} (FastAPI) e o banco de dados (PostgreSQL), e
Vercel para a entrega do \textit{front-end} (React). A meta de disponibilidade
é de 99,5\% de \textit{uptime} mensal, garantindo que interrupções no
\textit{software} não paralisem as atividades da oficina. \\
\hline
RNF04 & Portabilidade &
A interface deve seguir os princípios de \textit{Responsive Web Design},
assegurando pleno funcionamento em dispositivos móveis. Esta exigência é
derivada diretamente da fase de observação direta: mecânicos necessitam
registrar avarias e fotografias ao lado do veículo, sem deslocamento até uma
estação de trabalho fixa. \\
\hline
RNF05 & Conformidade Legal (LGPD) &
O sistema deve observar a Lei Geral de Proteção de Dados (Lei
n\textordmasculine~13.709/2018), em especial o Art.~18, que garante ao titular
os direitos de: (i) acesso aos dados pessoais armazenados; (ii) portabilidade
dos dados mediante solicitação; (iii) correção de dados incompletos ou
desatualizados; e (iv) exclusão dos dados quando encerrado o relacionamento
contratual (\textit{direito ao esquecimento}). O sistema deve registrar o
consentimento do titular no ato do cadastro. Este é um \textbf{requisito de
domínio} derivado do arcabouço jurídico brasileiro
(ver Subseção~\ref{subsubsec:dominio}). \\
\hline
RNF06 & Arquitetura de Comunicação &
O \textit{back-end} (FastAPI) deve expor uma API RESTful com documentação
gerada automaticamente no padrão OpenAPI/Swagger, facilitando futuras
integrações com fornecedores de peças, aplicativos móveis de terceiros e
o portal do cliente (RF07). \\

\end{longtable}



% ------------------------------------------------------------
\subsubsection{4.5.2.3 Requisitos de Domínio}
\label{subsubsec:dominio}
% ------------------------------------------------------------

Conforme \citeonline{sommerville2011}, requisitos de domínio são restrições e
características derivadas do domínio de aplicação do sistema, e não
diretamente das necessidades expressas pelos usuários. Eles refletem
convenções, regulamentações e práticas consolidadas no setor de atuação do
\textit{software}. No contexto do Torque Gestão, dois requisitos de domínio
foram identificados:

\begin{enumerate}
    \item \textbf{Formato de Placa Mercosul (RF02):} A legislação brasileira
    determinou, por meio da Resolução CONTRAN n\textordmasculine~729/2018, a
    adoção progressiva do padrão alfanumérico Mercosul (e.g., ABC1D23) para
    novos emplacamentos. O sistema deve validar placas tanto no formato
    antigo (AAA-9999) quanto no novo padrão, garantindo compatibilidade com
    o parque veicular heterogêneo atendido pelas oficinas.

    \item \textbf{Conformidade com a LGPD (RNF05):} A Lei Geral de Proteção
    de Dados impõe obrigações específicas a qualquer sistema que armazene
    dados pessoais de clientes, independentemente da vontade dos
    \textit{stakeholders}. O não atendimento a este requisito expõe o
    estabelecimento a sanções administrativas da ANPD (Autoridade Nacional
    de Proteção de Dados).
\end{enumerate}

% ------------------------------------------------------------
\subsection{4.5.3 Validação de Requisitos}
\label{subsec:validacao}
% ------------------------------------------------------------

A validação de requisitos é a etapa destinada a verificar se a especificação
produzida representa, de fato, o sistema que o cliente deseja construir
\cite{sommerville2011}. Diferentemente da verificação — que pergunta
\textit{"o sistema foi construído corretamente?"} —, a validação pergunta
\textit{"o sistema correto foi especificado?"}. \citeonline{sommerville2011}
define cinco critérios fundamentais que um conjunto de requisitos deve
satisfazer: validade, consistência, completude, realismo e verificabilidade.
Todos foram aplicados neste projeto, conforme descrito a seguir.

\subsubsection{4.5.3.1 Validade}

Verificou-se se as funções especificadas correspondem às necessidades reais dos
usuários identificadas na fase de elicitação. Para tanto, os cenários e
\textit{User Stories} construídos na
Subseção~\ref{subsec:elicitacao} foram utilizados como instrumento de
rastreabilidade: cada RF foi confrontado com pelo menos uma narrativa de uso,
garantindo que nenhum requisito foi inserido por suposição da equipe sem
respaldo nas entrevistas ou na observação direta.

\subsubsection{4.5.3.2 Consistência}

Os requisitos foram revisados em conjunto para detectar contradições internas.
Um potencial conflito identificado foi entre RNF02 (desempenho — consultas em
até 2~segundos) e RF03 (OS com separação detalhada de itens, operação de maior
custo computacional). A análise concluiu que o uso de índices relacionais no
PostgreSQL e a paginação de resultados no \textit{back-end} FastAPI são
suficientes para que ambos os requisitos sejam satisfeitos simultaneamente,
não havendo, portanto, inconsistência irreconciliável.

\subsubsection{4.5.3.3 Completude}

Verificou-se se todos os fluxos de trabalho identificados na etnografia estavam
cobertos pela especificação. A revisão cruzada entre os cenários de uso
(Subseção~\ref{subsec:elicitacao}) e as tabelas de requisitos funcionais
confirmou que os processos de cadastro de clientes e veículos, abertura e
acompanhamento de OS, controle de acesso por perfil, portal de acompanhamento
do cliente (RF07) e conformidade legal foram integralmente contemplados.
O escopo do sistema delimita-se à gestão operacional da oficina; funcionalidades
de natureza fiscal e tributária — como emissão de notas fiscais e controle
contábil — estão fora do escopo do Torque Gestão e devem ser endereçadas por
sistemas especializados externos. As funcionalidades identificadas como
desejáveis, mas inviáveis dentro do escopo do MVP (\textit{Minimum Viable
Product}) — como a integração com fornecedores de peças via API — foram
registradas como requisitos postergados para versões futuras, garantindo a
completude do documento sem comprometer a viabilidade da entrega.

\subsubsection{4.5.3.4 Realismo}

Avaliou-se se o orçamento e o cronograma do PAC Extensionista permitiam a
implementação de todos os requisitos especificados. Funcionalidades de alta
complexidade foram postergadas para versões futuras com o objetivo de garantir
a entrega de um MVP estável e funcional dentro do prazo acadêmico estabelecido.
Esta análise foi conduzida em reuniões com o professor orientador Claudinei
Dias, que validou o escopo frente à capacidade técnica e ao tempo disponível da
equipe.

\subsubsection{4.5.3.5 Verificabilidade}

Cada requisito foi redigido de forma que pudesse ser submetido a testes
objetivos e mensuráveis. O RNF02, por exemplo, pode ser verificado por
ferramentas de monitoramento de latência de API, como o \textit{Postman} ou o
painel de métricas do Render. O RF06 pode ser testado por meio de cenários de
controle de acesso automatizados, tentando acessar rotas administrativas com
credenciais de perfil mecânico ou cliente e verificando a resposta HTTP~403
esperada. O RF07 pode ser validado por meio de testes de integração que
confirmem o isolamento de dados entre clientes distintos e a correta exibição
do status de OS em tempo real. A adoção deste critério desde a escrita dos
requisitos assegura que a fase de testes do sistema terá insumos claros para a
construção dos casos de teste.

\bigskip
\noindent
O processo de Engenharia de Requisitos aqui documentado garante que o Torque
Gestão não seja apenas um \textit{software} funcionalmente correto, mas uma
solução estratégica, juridicamente adequada e operacionalmente confiável para
o mercado automotivo de Joinville. A rastreabilidade entre as necessidades
identificadas na elicitação, os requisitos especificados e os critérios de
validação aplicados constitui a base sobre a qual as fases subsequentes de
projeto, implementação e testes serão conduzidas.
 
% 5 AVALIAÇÃO DO PROJETO
\newpage
\phantomsection
\label{sec:avaliacao}
\section*{5 AVALIA\c{C}\~AO DO PROJETO PELO P\'UBLICO BENEFICIADO}
 
% 6 CONSIDERAÇÕES FINAIS
\newpage
\phantomsection
\label{sec:consideracoes}
\section*{6 CONSIDERA\c{C}\~OES FINAIS -- AUTOAVALIA\c{C}\~AO DO PAC EXTENSIONISTA}
 
% REFERÊNCIAS EM BRANCO
\newpage
\phantomsection
\label{sec:referencias}
\section*{REFER\^ENCIAS}
\bibliography{referencias}
 
% APÊNDICES EM BRANCO
\newpage
\phantomsection
\label{sec:apendices}
\section*{AP\^ENDICES}
 
% ANEXOS EM BRANCO
\newpage
\phantomsection
\label{sec:anexos}
\section*{ANEXOS}
 
\end{document}